\begin{figure}
    \centering
    \begin{minipage}{0.49 \linewidth}
    \includegraphics[width=\linewidth]{images/adaptivecifarlr1.pdf}
    \end{minipage} \hfill
    \begin{minipage}{0.49 \linewidth}
    \includegraphics[width=\linewidth]{images/adaptivecifarlr3.pdf}
    \end{minipage}
    \caption{CIFAR-10 training loss with fixed and adaptive $\alpha$. The adaptive $\alpha$ is clipped between $[\alpha_{low}, 1]$. (Left) Adam learning rate = 0.001.  (Right) Adam learning rate = 0.003.} 
    \label{fig:adaptivealpha}
\end{figure}

\subsection{Selecting the Slow Weights Step Size}

The step size in the direction ($\theta_{t,k} - \theta_{t,0}$) is controlled by $\alpha$. By taking a quadratic approximation of the loss, we present a principled way of selecting $\alpha$. 

\begin{proposition}[Optimal slow weights step size]
\label{prop:optimal-step-size}
For a quadratic loss function $L(x) = \frac{1}{2} x^T A x - b^T x$, 
the step size $\alpha^*$ that minimizes the loss for two points $\theta_{t,0}$ and $\theta_{t,k}$ is given by:
\[ \alpha^* = \arg \min_{\alpha} L(\theta_{t,0} + \alpha (\theta_{t,k} - \theta_{t,0})) = \frac{(\theta_{t,0} - \theta^*)^T A(\theta_{t,0} - \theta_{t, k})}{(\theta_{t,0}  - \theta_{t,k})^T A (\theta_{t,0} - \theta_{t,k})}  \]
% \[ \alpha^* = \arg \min_{\alpha} L(x_n + \alpha (x_{n+1} - x_n)) = \frac{(x_n - x^*)^T A(x_n - x_{n+1})}{(x_n - x_{n+1})^T A(x_n - x_{n+1})}  \]
where $\theta^* = A^{-1}b$ minimizes the loss.
\end{proposition}
% \begin{proof}
% Straightforward, set the derivative with respect to $\alpha$ to 0. Details in appendix.
% \end{proof}
%By Taylor's Theorem, if $A$ represents the true Hessian of the loss, this choice of $\alpha^*$ is optimal up to third order terms. \MZ{Not sure Taylor Theorem is accurate, but we should justify using a quadratic approx.} 
Proof is in the appendix. Using quadratic approximations for the curvature, which is typical in second order optimization \citep{duchi2011adaptive, kingma2014adam, martens2015optimizing}, we can derive an estimate for the optimal $\alpha$ more generally. The full Hessian is typically intractable so we instead use aforementioned approximations, such as the diagonal approximation to the empirical Fisher used by the Adam optimizer \citep{kingma2014adam}. This approximation works well in our numerical experiments if we clip the magnitude of the step size. At each slow weight update, we compute:

% \[ \hat{\alpha}^* = \text{clip}( \frac{(x_n - (x_{n+1} - \hat{A}^{-1} \nabla x_{n+1} )^T \hat{A} (x_n - x_{n+1})}{(x_n - x_{n+1})^T \hat{A} (x_n - x_{n+1})},  \alpha_{\text{low}}, 1)\]

\[ \hat{\alpha}^* = \text{clip}( \frac{(\theta_{t,0} - (\theta_{t,k}  - \hat{A}^{-1} \hat{\nabla} L(\theta_{t,k}) )^T \hat{A} (\theta_{t,0}  - \theta_{t,k})}{(\theta_{t,0}  - \theta_{t,k})^T \hat{A} (\theta_{t,0} - \theta_{t,k})},  \alpha_{\text{low}}, 1)\]


% \MZ{perhaps we should have an algorithm box for this in the appendix?}

 where $\hat{A}$ is the empirical Fisher approximation and $\theta_{t,k}  - \hat{A}^{-1} \hat{\nabla} L(\theta_{t,k})$ approximates the optimum $\theta^*$. We prove Proposition \ref{prop:optimal-step-size} and elaborate on assumptions in the appendix \ref{app:optimal-slow}. Setting $\alpha_{\text{low}} > 0$ improves the stability of our algorithm. We evaluate the performance of this adaptive scheme versus a fixed scheme and standard Adam on a ResNet-18 trained on CIFAR-10 with two different learning rates and show the results in Figure \ref{fig:adaptivealpha}. Additional hyperparameter details are given in appendix \ref{app:experiments}. Both the fixed and adaptive Lookahead offer improved convergence.

In practice, a fixed choice of $\alpha$ offers similar convergence benefits and tends to generalize better. Fixing $\alpha$ avoids the need to maintain an estimate of the empirical Fisher, which incurs a memory and computational cost when the inner optimizer does not maintain such an estimate e.g. SGD. We thus use a fixed $\alpha$ for the rest of our deep learning experiments. % to be elaborated on